动态实验中，28/30个订单流的四种机制均可执行，相对禁合作净收益差均值为+58.9（16/12/0），相对无参与底线净收益差均值为-26.3（6/14/8），相对顺序插入基线净收益差均值为+193.2（25/3/0）；保留2个受控不可执行单元，另有2个阶段超过下一触发间隔，故只支持批量滚动决策解释。固定配送方案在28个电网日的充电择时复算显示，六个单元的充电排放降幅为2.78\%--3.28\%。这些结果共同说明动态协同的收益、参与保障、可执行性、计算时限和充电减排必须分项判断，不能压缩成单一优化目标。
